\documentclass[11pt]{article}
\usepackage{amsmath,amssymb,amsthm}
\usepackage{hyperref}

\newtheorem{theorem}{Theorem}
\newtheorem{conjecture}[theorem]{Conjecture}
\newtheorem{remark}[theorem]{Remark}
\newtheorem{observation}[theorem]{Observation}

\newcommand{\CPP}{\operatorname{CPP}}
\newcommand{\GK}{\operatorname{GK}}
\newcommand{\QSym}{\operatorname{QSym}}
\newcommand{\Sym}{\operatorname{Sym}}
\newcommand{\qpoch}[2]{(#1;#2)_\infty}
\newcommand{\qpochn}[3]{(#1;#2)_{#3}}
\newcommand{\sln}{\widehat{\mathfrak{sl}}_n}

\title{Literature Review: Three Threads Toward Conjecture~2.7\\
{\large Tingley, Tight Partitions, and Quasisymmetric Functions}}
\author{Robin Langer}
\date{February 2026}

\begin{document}
\maketitle

\begin{abstract}
We survey three approaches to Warnaar's Conjecture~2.7
(positivity of $Q_{n,c}(q)$ for cylindric plane partitions)
and assess what each contributes and where each hits its limit.
The three threads are: (1)~Tingley's crystal base interpretation
of cylindric partitions, (2)~Kanade--Russell's tight cylindric
partitions and functional equations, and (3)~our own
quasisymmetric function refinement and $L$-positivity results.
All three illuminate different facets of the conjecture but
each encounters the same obstacle: the $[z^n]$ extraction
that isolates CPPs with a fixed maximum entry.
\end{abstract}

%% ========================================================================
\section{Thread 1: Tingley --- Crystal bases and cylindric partitions}
\label{sec:tingley}
%% ========================================================================

\subsection{Main result}

Tingley~\cite{tingley-2008} proved that cylindric plane partitions
are crystal basis elements of $\sln$. He establishes three
equivalent combinatorial models for the crystal graph:

\begin{enumerate}
  \item \textbf{Hayashi realization}: partitions with crystal
    operators $e_i$, $f_i$ acting by adding/removing boxes
    colored by residue mod~$n$.
  \item \textbf{Abacus model}: beads on an $n$-strand abacus,
    with crystal operators sliding beads between strands.
  \item \textbf{Kyoto path model}: semi-infinite tensor products
    of elements from a perfect Kirillov--Reshetikhin crystal
    $B_{r,s}$.
\end{enumerate}

\begin{theorem}[Tingley, Theorem~3.14]
The colored directed graph formed by tight, descending abacus
configurations is a realization of the crystal graph for the
$\sln$ representation $V(\Lambda)$.
\end{theorem}

The central application: Tingley gives a
\textbf{representation-theoretic proof of Borodin's product
formula}~\cite{borodin-2007}. The generating function
$\GK_c(1,q)$ is identified with the principally specialized
character of an integrable highest-weight $\sln$-module,
and characters of integrable modules factor by the
Weyl--Kac character formula.

\subsection{The underlying chain}

The connection runs through the corner transfer matrix (CTM)
method of Baxter. In a solvable lattice model satisfying the
Yang--Baxter equations, the CTM eigenstates at $q = 0$ form the
crystal basis $B(\Lambda)$. A \emph{path} through this crystal
is a sequence $(b_1, b_2, \ldots, b_L) \in B_{r,s}^{\otimes L}$.
When the path is periodic (ground-state boundary condition),
it becomes a cylindric partition. The energy function
$H : B_1 \times B_2 \to \mathbb{Z}$ tracks the $q$-weight:
\[
  X(B, \Lambda; q)
  = \sum_{\text{paths } p} q^{\operatorname{energy}(p)}
  = \GK_c(1, q).
\]
This one-dimensional configuration sum equals the CPP generating
function, which factors into Borodin's product.

\subsection{The bijection chain}

There is a chain of weight-preserving bijections:
\[
  \text{Cylindric partitions}
  \xleftrightarrow{\text{Tingley}}
  \text{Crystal paths}
  \xleftrightarrow{\text{Kyoto}}
  \text{Semi-infinite tensor products}
  \xleftrightarrow{\text{OSS}}
  \text{Rigged configurations,}
\]
where OSS refers to Okado--Schilling--Shimozono~\cite{oss-2001}.
Each arrow preserves the generating function weight
(energy $\leftrightarrow$ size $\leftrightarrow$ cocharge).

\subsection{Relevance to Conjecture~2.7}

If $Q_{n,c}(q)$ could be interpreted as a weight multiplicity
or branching coefficient inside the crystal, positivity would be
automatic---crystals count things. The obstacle is that the
$[z^n]$ extraction (fixing the maximum entry of a CPP) does not
have a known crystal interpretation. The maximum entry
$n(\pi) = \max\{\nu^{(i)}_j\}$ is not a crystal-theoretic
statistic: it depends on the absolute values of entries, not on
their relative positions or the action of the operators $e_i$,
$f_i$.

\begin{remark}
Tingley's crystal structure lives on the \emph{full} set of CPPs
(all max entries). Restricting to $\max = n$ breaks the crystal
into fragments that do not form sub-crystals in general.
This parallels the failure of quasi-symmetry at fixed max entry
(Thread~3).
\end{remark}

\subsection{Key references}

\begin{itemize}
  \item Tingley~\cite{tingley-2008}: the central bridge paper.
  \item Kang--Kashiwara--Misra~\cite{kkm-1992}: perfect crystals
    and the path realization theorem.
  \item Okado--Schilling--Shimozono~\cite{oss-2001}: path-to-rigged-configuration bijection.
  \item Nakayashiki--Yamada~\cite{ny-1995}: Kostka polynomials and
    energy functions in solvable lattice models.
\end{itemize}

%% ========================================================================
\section{Thread 2: Kanade--Russell --- Tight cylindric partitions}
\label{sec:tight}
%% ========================================================================

\subsection{Main result}

Kanade and Russell~\cite{kanade-russell-2023} define
\emph{tight} cylindric partitions, where consecutive stack
heights differ by at most~1 (a stricter condition than ordinary
interlacing). They prove:

\begin{theorem}[Kanade--Russell, Theorem~4]
For a profile $c = (c_1, \ldots, c_r)$ with
$m = c_1 + \cdots + c_r$, the univariate generating function
for tight CPPs is
\[
  T_c(1,q)
  = \Pr\bigl(L_{\sln[r]}(\Lambda(c))\bigr)
  = \frac{(q^r; q^r)_\infty\, (q^m; q^m)_\infty^{r-1}}
         {(q)_\infty^r}
    \prod_{1 \le i < j \le r}
    \theta\bigl(q^{j - i + c_i + \cdots + c_{j-1}};\, q^m\bigr).
\]
\end{theorem}

The proof uses \textbf{Corteel--Welsh functional equations}:
a system of equations for the bivariate generating function
$T_c(z, q)$, solved via inclusion-exclusion and induction on
$z$-exponents. This is a purely analytic approach---no
representation theory, no crystal bases.

\subsection{Additional results}

\begin{itemize}
  \item \textbf{Theorem~10}: Corteel--Welsh functional equation
    analogue for $r$-rowed tight CPPs.
  \item \textbf{Theorem~13}: Bivariate closed forms for
    2-rowed tight CPPs. For profile $(c_1, c_2)$, the generating
    functions reduce to $(-q;q)_\infty$ times Gordon--Andrews or
    Andrews--Bressoud identities.
  \item \textbf{Theorem~14}: Diamond relations reformulation.
  \item The 3-row case remains open (only conjectural bivariate
    formulas for small profiles).
\end{itemize}

\subsection{Relevance to Conjecture~2.7}

The Corteel--Welsh functional equations are powerful for
computing generating functions and proving product formulas,
but they do \emph{not} directly address positivity. The explicit
double-sum formulas for $[y^n] G_c(y,q)$ (Theorem~3.2 of
Corteel--Welsh~\cite{corteel-welsh-2019}) involve terms with
alternating signs for profiles like $(3,0,1)$ and $(2,2,0)$.
Positivity of $Q_{n,c}(q)$ requires delicate cancellation that
the functional equations cannot explain.

The functional equations are a computational oracle, not a
conceptual explanation. They tell us \emph{what} $Q_{n,c}(q)$
equals, but not \emph{why} its coefficients should be positive.

\begin{remark}
The tight CPP product formula is structurally parallel to
Borodin's identity (both are principally specialized
$\sln[r]$ characters). This suggests that an alternative proof
via crystal bases (extending Tingley's approach to tight CPPs)
should be possible and would deepen the representation-theoretic
picture. It could also reveal whether analogous $Q$-polynomials
exist for the tight case.
\end{remark}

%% ========================================================================
\section{Thread 3: Quasisymmetric functions and $L$-positivity}
\label{sec:qsym}
%% ========================================================================

\subsection{Main results}

In~\cite{langer-qsym}, we showed that the CPP frequency
generating function
\[
  \widetilde{\GK}_c(\mathbf{z})
  = \sum_{\pi \in \CPP(c)}
    z_1^{\#_1(\pi)}\, z_2^{\#_2(\pi)} \cdots
\]
(where $\#_k(\pi)$ counts parts equal to~$k$) lies in the ring
$\QSym$ of quasisymmetric functions. This follows from Stanley's
$P$-partition theory: the interlacing conditions depend only on
the relative order of entries, not their absolute values.

\begin{observation}[Verified: 55 profiles, $d \leq 5$]
$\widetilde{\GK}_c$ is $L$-positive:
$\widetilde{\GK}_c = \sum_\alpha b_\alpha(c)\, L_\alpha$
with $b_\alpha \in \mathbb{Z}_{\geq 0}$.
\end{observation}

The stabilized $L$-coefficients follow $b^{\ell(\alpha)}$
patterns ($b = 2$ for $d = 2$, $b = 5$ for $d = 4$), suggesting
a product structure. By Tewari--van Willigenburg~\cite{tvw-2015},
$L$-positivity is equivalent to the existence of a $0$-Hecke
module whose character is $\widetilde{\GK}_c$.

\begin{conjecture}[Descent statistic]
There exists a natural descent-type statistic on cylindric
partitions such that
$\widetilde{\GK}_c = \sum_\pi L_{\alpha(\pi)}$,
where $\alpha(\pi)$ is the descent composition of~$\pi$.
\end{conjecture}

\subsection{The dead end}

$L$-positivity does \textbf{not} extend to the $Q$-polynomials.
Two obstacles:

\begin{enumerate}
  \item \textbf{Max-entry constraint breaks quasi-symmetry.}
    The individual slices
    $\tilde{g}_n = \sum_{\max(\pi) = n} z_1^{\#_1} \cdots$
    are \emph{not} quasi-symmetric: fixing $\max = n$ privileges
    the value~$n$, breaking the relabelling argument.
  \item \textbf{Cross-frequency cancellation.}
    The alternating-sign factor $(zq;q)_\infty$ introduces
    cancellation that crosses between different frequency vectors.
    Computationally, $\sim$3--6\% of the cancellation in $Q_n$ is
    cross-frequency: the net coefficient at individual frequency
    vectors is negative, and only the aggregation across all
    compositions of the same weight (under $z_k \mapsto q^k$)
    yields the positive $Q_n$ coefficient.
\end{enumerate}

The positivity of $Q_{n,c}(q)$ is an
\textbf{intrinsically one-variable phenomenon}: the
specialization $z_k \mapsto q^k$ is essential and cannot be
avoided within $\QSym$.

\subsection{Connection to Lam--Pylyavskyy Conjecture~5.5}

Conjecture~5.5 of~\cite{lam-pylyavskyy-2006} characterizes when
products $L_\alpha \cdot L_\beta$ are maximal in the
$L$-positivity order (iff the pair $\{\alpha, \beta\}$ is stable
under cell transfer). This is relevant because the stabilization
pattern of $\widetilde{\GK}_c$ suggests a product structure in
the $L$-basis, and Conjecture~5.5 would constrain its form.
However, this enriches the understanding of $\widetilde{\GK}_c$
itself (the unconstrained generating function) without bridging
the gap to $Q_{n,c}$ positivity.

%% ========================================================================
\section{Synthesis}
\label{sec:synthesis}
%% ========================================================================

The three threads illuminate different facets of the same problem:

\medskip
\begin{center}
\begin{tabular}{l|l|l}
  \textbf{Thread} & \textbf{What it gives} & \textbf{Where it stops} \\
  \hline
  Tingley (crystals) & \emph{Why} Borodin's product holds &
    $[z^n]$ extraction has no crystal \\
  & (characters of $\sln$ modules) & interpretation \\
  \hline
  Kanade--Russell & \emph{What} $G_c(y,q)$ equals &
    Formulas have mixed signs; \\
  (functional eqns) & (explicit double sums) & no positivity mechanism \\
  \hline
  QSym / $L$-positivity & \emph{Structure} of $\widetilde{\GK}_c$ &
    Positivity is one-variable; \\
  & ($L$-positive, $0$-Hecke module) & cross-frequency cancellation
\end{tabular}
\end{center}

\medskip

All three approaches converge on the same obstacle: the $[z^n]$
extraction that isolates CPPs with a specific maximum entry.
In the crystal picture, fixing max entry breaks the crystal into
non-sub-crystal fragments. In the QSym picture, it breaks
quasi-symmetry. In the functional equation picture, it
introduces alternating signs.

\subsection{What would a proof need?}

A proof of Conjecture~2.7 likely requires one of:
\begin{enumerate}
  \item A \textbf{combinatorial interpretation} of $Q_{n,c}(q)$
    as a counting function---finding what $Q_{n,c}$ counts.
  \item A \textbf{crystal-theoretic interpretation} of the
    $[z^n]$ extraction---expressing $Q_{n,c}$ as a weight
    multiplicity or branching coefficient.
  \item A \textbf{new positivity mechanism} that works in one
    variable---perhaps a positivity-preserving transformation
    (in the spirit of the Borwein program) or a connection to
    Kazhdan--Lusztig theory.
\end{enumerate}

The most promising direction may be~(2): if the
``$\max = n$ slice'' of the crystal corresponds to some natural
representation-theoretic object (a Demazure module? a level
restriction? a truncated path space?), then positivity would
follow from the general theory. This is the direction suggested
by the parallel between Borodin's identity and the
Kanade--Russell product formula---both are $\sln[r]$ characters,
and the extraction of $Q_n$ should have a meaning in that world.

%% ========================================================================
%% References
%% ========================================================================

\begin{thebibliography}{99}

\bibitem{borodin-2007}
A.~Borodin,
Periodic Schur process and cylindric partitions,
\emph{Duke Math.\ J.}\ \textbf{140}(3) (2007), 391--468.
arXiv:math/0601019.

\bibitem{corteel-welsh-2019}
S.~Corteel and T.~Welsh,
The $A_2$ Rogers--Ramanujan identities revisited,
arXiv:1905.08343, 2019.

\bibitem{kanade-russell-2023}
S.~Kanade and M.~C.~Russell,
Tight cylindric partitions,
arXiv:2507.02435, 2023.

\bibitem{kkm-1992}
S.-J.~Kang, M.~Kashiwara, and K.~C.~Misra,
Crystal bases of Verma modules for quantum affine Lie algebras,
\emph{Compositio Math.}\ \textbf{92}(3) (1994), 299--325.

\bibitem{lam-pylyavskyy-2006}
T.~Lam and P.~Pylyavskyy,
$P$-partition products and fundamental quasi-symmetric function
positivity,
\emph{Adv.\ in Appl.\ Math.}\ \textbf{40}(3) (2008), 331--354.

\bibitem{langer-qsym}
R.~Langer,
Quasi-symmetric refinement of the cylindric partition
generating function,
In preparation, 2025.

\bibitem{ny-1995}
A.~Nakayashiki and Y.~Yamada,
Kostka polynomials and energy functions in solvable lattice
models,
\emph{Selecta Math.\ (N.S.)}\ \textbf{3}(4) (1997), 547--599.
arXiv:q-alg/9512027.

\bibitem{oss-2001}
M.~Okado, A.~Schilling, and M.~Shimozono,
Virtual crystals and fermionic formulas of type
$D^{(2)}_{n+1}$, $A^{(2)}_{2n}$, and $C^{(1)}_n$,
\emph{Represent.\ Theory}\ \textbf{7} (2003), 101--163.
arXiv:math/0105017.

\bibitem{tingley-2008}
P.~Tingley,
Three combinatorial models for $\widehat{\mathfrak{sl}}_n$
crystals, with applications to cylindric plane partitions,
\emph{Int.\ Math.\ Res.\ Not.}\ 2008(2), Art.~ID rnm143.
arXiv:math/0702062.

\bibitem{tvw-2015}
V.~V.~Tewari and S.~J.~van Willigenburg,
Modules of the $0$-Hecke algebra and quasisymmetric Schur
functions,
\emph{Adv.\ Math.}\ \textbf{285} (2015), 1025--1065.
arXiv:1403.1527.

\end{thebibliography}

\end{document}
